\chapter{Complexity — Layered P Poly}
%%% generated by the blueprint pipeline from TCSlib.Complexity.CircuitComplexity.LayeredPPoly
\section{Overview}

Declarations of the Lean module \texttt{TCSlib.Complexity.CircuitComplexity.LayeredPPoly}.

\section{Declarations}

\begin{definition}[Layered circuit family]
\lean{BoolCircuit.LayeredCircuitFamily}
\statementsource{lmn-constant-depth-fourier}{pdf-63759bd909d5-p004-b007}
\label{BoolCircuit.LayeredCircuitFamily}
\leanok
\uses{BoolCircuit.LayeredCircuit, BoolCircuit.LayeredCircuit.Finite}
A layered circuit family over the binary alphabet consists of one layered circuit for
each input length $n$, with inputs indexed by the $n$ positions and output in the
one-point set. Each circuit in the family is finite.
\end{definition}

\begin{definition}[Acceptance by a circuit family]
\lean{BoolCircuit.LayeredCircuitFamily.Accepts}
\label{BoolCircuit.LayeredCircuitFamily.Accepts}
\leanok
\uses{BoolCircuit.LayeredCircuit.eval₁}
Let $C$ be a family of layered circuits indexed by input length. A Boolean word $w$ is
accepted by $C$ exactly when the circuit indexed by the length of $w$, evaluated on the
bits of $w$, has output $1$.
\end{definition}

\begin{definition}[Language of a circuit family]
\lean{BoolCircuit.LayeredCircuitFamily.language}
\label{BoolCircuit.LayeredCircuitFamily.language}
\leanok
\uses{BoolCircuit.LayeredCircuitFamily.Accepts}
For a family $C$ of layered circuits, define its language to be the set of Boolean words
accepted by $C$.
\end{definition}

\begin{theorem}[Membership in the circuit family language]
\lean{BoolCircuit.LayeredCircuitFamily.mem_language_iff}
\label{BoolCircuit.LayeredCircuitFamily.mem_language_iff}
\leanok
\uses{BoolCircuit.LayeredCircuit.eval₁, BoolCircuit.LayeredCircuitFamily.language}
\difficulty{1}
Let $C$ be a family of layered circuits and let $w$ be a finite Boolean word. Then $w$
belongs to the language accepted by $C$ if and only if the circuit in the family for
inputs of length $|w|$, evaluated on the input obtained by encoding each bit of $w$ as
its corresponding element of $\mathrm{Fin}\,2$, outputs $1$.
\end{theorem}

\begin{definition}[Gate restriction for a circuit family]
\lean{BoolCircuit.LayeredCircuitFamily.OnlyUsesGates}
\statementsource{lmn-constant-depth-fourier}{pdf-63759bd909d5-p004-b007}
\label{BoolCircuit.LayeredCircuitFamily.OnlyUsesGates}
\leanok
\uses{BoolCircuit.GateOp, BoolCircuit.LayeredCircuit.onlyUsesGates}
Let $S$ be a set of gate operations on the binary alphabet. A layered circuit family
uses only gates from $S$ if, for every input length, its circuit uses only gate
operations belonging to $S$.
\end{definition}

\begin{definition}[Polynomial size circuit family]
\lean{BoolCircuit.LayeredCircuitFamily.IsPolySize}
\statementsource{lmn-constant-depth-fourier}{pdf-63759bd909d5-p004-b007}
\statementsource{mit-18405-razborov-smolensky}{
  pdf-d3f464808105-p001-b005,
  pdf-d3f464808105-p002-b001
}
\label{BoolCircuit.LayeredCircuitFamily.IsPolySize}
\leanok
\uses{BoolCircuit.LayeredCircuit.size}
A layered circuit family $C$ has polynomial size if there are natural numbers $a$ and
$k$ such that, for every input length $n$, the size of its circuit at length $n$ is at
most $a(n+1)^k$.
\end{definition}

\begin{definition}[Layered circuit size class]
\lean{Language.InLayeredSIZE}
\label{Language.InLayeredSIZE}
\leanok
\uses{BoolCircuit.LayeredCircuit.size, BoolCircuit.LayeredCircuitFamily, BoolCircuit.LayeredCircuitFamily.OnlyUsesGates, BoolCircuit.LayeredCircuitFamily.language, BoolCircuit.stdGateOps}
A language $L$ belongs to layered size $T$ if there is a family $C$ of layered Boolean
circuits such that every circuit uses only the standard gate operations, the circuit on
$n$ input bits has size at most $T(n)$ for every natural number $n$, and the language
recognized by $C$ is exactly $L$.
\end{definition}

\begin{definition}[Languages decidable in layered polynomial size]
\lean{Language.InLayeredPPoly}
\statementsource{lmn-constant-depth-fourier}{pdf-63759bd909d5-p004-b007}
\label{Language.InLayeredPPoly}
\leanok
\uses{Language.InLayeredSIZE}
A language $L$ over the Boolean alphabet is decidable by layered circuits of polynomial
size if there are natural numbers $a$ and $k$ such that $L$ has layered circuit
complexity at most the function $n\mapsto a(n+1)^k$.
\end{definition}

\begin{theorem}[Polynomial-size layered circuit characterization]
\lean{Language.inLayeredPPoly_iff}
\label{Language.inLayeredPPoly_iff}
\leanok
\uses{BoolCircuit.LayeredCircuitFamily, BoolCircuit.LayeredCircuitFamily.IsPolySize, BoolCircuit.LayeredCircuitFamily.OnlyUsesGates, BoolCircuit.LayeredCircuitFamily.language, BoolCircuit.stdGateOps, Language.InLayeredPPoly}
\difficulty{4}
A Boolean language $L$ is recognized by a polynomial-size family of layered circuits
using only identity, NOT, and unbounded fan-in AND gates if and only if there is such a
family whose circuit sizes are bounded by a polynomial in the input length and whose
recognized language is exactly $L$.
\end{theorem}

\begin{definition}[Class of layered polynomial size languages]
\lean{BoolCircuit.LayeredPPoly}
\statementsource{lmn-constant-depth-fourier}{pdf-63759bd909d5-p004-b007}
\label{BoolCircuit.LayeredPPoly}
\leanok
\uses{Language.InLayeredPPoly}
The class of layered polynomial size languages consists of all languages over the
Boolean alphabet that are decidable by layered circuits of polynomial size.
\end{definition}

\begin{theorem}[Characterization of layered polynomial size]
\lean{BoolCircuit.mem_LayeredPPoly_iff}
\label{BoolCircuit.mem_LayeredPPoly_iff}
\leanok
\uses{BoolCircuit.LayeredPPoly, Language.InLayeredPPoly}
\difficulty{1}
For every language $L$ over the Boolean alphabet, $L$ belongs to the class of languages
with polynomial-size layered circuits using the specified standard gate operations if
and only if there are natural numbers $a$ and $k$ and such a circuit family recognizing
exactly $L$ whose size on inputs of length $n$ is at most $a(n+1)^k$.
\end{theorem}
